% =============================================================================
%  s05_seccion_cierre.tex — Sección 5: Conclusiones y Cierre
% =============================================================================

\section{Conclusiones y Cierre}

% ── Frame 5.1: El Transporte como Espacio Social ──────────────────────────────
%
%  Diseño:
%  Gancho narrativo al tope (alien quote centrada, dos reglas doradas).
%  Dos columnas simétricas: block verde (confirmaciones + paráfrasis espacial)
%  y alertblock rojo (retos pendientes + Segundo Espacio).
%  Pie: remate que conecta el 53% de mejora con el espacio vivido de millones.
%
\begin{frame}{El Transporte no es solo Infraestructura}
  % ── Gancho narrativo ─────────────────────────────────────────────────────────
  \begin{center}
    {\color{ColorAcento}\rule{0.88\textwidth}{0.6pt}}\\[0.07cm]
    {\scriptsize\itshape
      «Si una raza alienígena viera la Tierra desde el espacio, vería a dos razas
      conviviendo en el planeta; una que usa sus dos piernas para desplazarse
      y otra más dominante que corre a gran velocidad en cuatro ruedas.»}\\[0.03cm]
    {\tiny\color{gray!60}
      --- Inspirado en \citet{adams1979hitchhiker}}\\[0.06cm]
    {\color{ColorAcento}\rule{0.60\textwidth}{0.5pt}}
  \end{center}
  \vspace{0.06cm}
  \begin{columns}[T]
    % ── Lo que el modelo ya confirma ─────────────────────────────────────────
    \begin{column}{0.50\textwidth}
      \begin{block}{\scriptsize\bfseries Lo que el modelo ya confirma}
        {\scriptsize
          {\color{green!60!black}$\surd$}
            Cobertura sur: 11.75\,\% en Milpa Alta\\[0.03cm]
          {\color{green!60!black}$\surd$}
            $\bar{DF}=1.545$ — rodeo estructural, no accidental\\[0.03cm]
          {\color{green!60!black}$\surd$}
            $FC_H=130$ — sin alternativa en periferia\\[0.03cm]
          {\color{green!60!black}$\surd$}
            11 sistemas validados dentro de rango\\[0.09cm]
          {\tiny\color{ColorPrincipal}
            \textbf{El transporte reproduce la segregación del espacio:}
            el \textbf{Tercer Espacio} de la periferia
            \citep{soja1996thirdspace} es espera y desvío —
            no de encuentro.
            El \textbf{Primer Espacio} material
            \citep{lefebvre2013produccion} llega tarde al sur.}}
      \end{block}
    \end{column}
    % ── Retos pendientes ─────────────────────────────────────────────────────
    \begin{column}{0.47\textwidth}
      \begin{alertblock}{\scriptsize\bfseries Retos pendientes}
        {\scriptsize
          Fase 3: $T$, $B(v)$, $\Delta E$ — pendiente\\[0.04cm]
          Vulnerabilidad crítica en nodos radiales\\
          sin caminos alternos transversales.\\[0.09cm]
          {\tiny\color{ColorPrincipal}
            El \textbf{Segundo Espacio} — planificadores
            y presupuestos — aún no incorpora
            la equidad territorial como restricción
            formal \citep{lefebvre2013produccion}.}%
        }
      \end{alertblock}
      \vspace{0.08cm}
      {\tiny\itshape\color{gray!60}
        La redundancia anillar reduce la vulnerabilidad
        en un $53\,\text{\%}$ — y con ella transforma
        el espacio vivido de millones.}
    \end{column}
  \end{columns}
\end{frame}

% ── Frame 5.2: La Higuera — Sylvia Plath (diapositiva completa) ───────────────
%
%  Frame completo centrado verticalmente.
%  Cita extensa de La campana de cristal: la imagen de los higos y la parálisis
%  de decisión — metáfora central del VFTModel (Vanishing Fig-Tree Model).
%  Remate en italica: el modelo mide cuántos higos quedan.
%
\begin{frame}{La Red que Desaparece en Silencio}
  \vspace{0.12cm}
  \begin{center}
    {\color{ColorAcento}\rule{0.72\textwidth}{0.9pt}}\\[0.22cm]
    {\small\itshape\color{ColorPrincipal}
      «Vi mi vida ramificarse ante mí como la higuera verde del cuento.\\[0.08cm]
      De la punta de cada rama, como un higo gordo y morado,\\
      me llamaba y me guiñaba un futuro maravilloso.\\[0.10cm]
      Un higo era un marido y un hogar feliz y unos hijos,\\
      y otro higo era una poeta famosa,\\
      y otro higo era una brillante profesora.\\[0.10cm]
      Yo quería cada uno de ellos,\\
      pero elegir uno significaba perder todos los demás.\\[0.10cm]
      Y mientras me sentaba allí, incapaz de decidirme,\\
      los higos empezaron a arrugarse y a tornarse negros,\\
      y, uno por uno, cayeron al suelo, a mis pies.»}\\[0.22cm]
    {\color{ColorAcento}\rule{0.72\textwidth}{0.9pt}}\\[0.14cm]
    {\scriptsize\color{gray!60}
      --- Sylvia Plath, \textit{La campana de cristal}
      \citep{plath2019campana}}\\[0.14cm]
    {\scriptsize\itshape\color{ColorPrincipal}
      Cada higo que cae es una zona periférica que pierde su conexión.\\
      El \texttt{VFTModel} mide cuántos higos quedan — y cuántos ya cayeron.}
    \vfill
  \end{center}
\end{frame}

% ── Frame 5.3: Estado Actual de la Tesis ──────────────────────────────────────
%
%  Tabla de 3 columnas: Cap · Contenido · Estado.
%  Misma paleta de colores y símbolos que la tabla de indicadores (frame 4.6).
%  Refleja fielmente el índice del documento de tesis (febrero 2026).
%
\begin{frame}{Estado Actual de la Tesis}
  \vspace{0.05cm}
  {\renewcommand{\arraystretch}{1.22}
  \begin{tabularx}{\textwidth}{@{\hspace{0.08cm}} c X c @{\hspace{0.08cm}}}
    \toprule
    {\tiny\textbf{Cap.}} &
    {\tiny\textbf{Contenido}} &
    {\tiny\textbf{Estado}} \\
    \midrule
    {\scriptsize 1–3} &
    {\tiny Introducción · Antecedentes y Marco Teórico · Definiciones} &
    {\tiny\color{green!60!black}$\surd$ Completo} \\[0.02cm]
    {\scriptsize 4} &
    {\tiny Capas de Código — \texttt{Apimetro} y \texttt{VFTModel} (código entregado)} &
    {\tiny\color{green!60!black}$\surd$ En revisión} \\
    \midrule
    {\scriptsize 5.1} &
    {\tiny Fase 1: Cobertura ($C$), Fuerza Capilar ($FC_H$),
           Factor de Desviación ($DF$)} &
    {\tiny\color{green!60!black}$\surd$ Calculado} \\[0.02cm]
    {\scriptsize 5.2} &
    {\tiny Fase 2: Fricción Vial ($CF$), Penalización Transbordo ($T_b$),
           Node Score} &
    {\tiny\color{orange!80!black}$\triangle$ Parcial} \\[0.02cm]
    {\scriptsize 5.3} &
    {\tiny Fase 3: Tiempo de Viaje ($T$), Intermediación ($B(v)$),
           Vulnerabilidad ($\Delta E$)} &
    {\tiny\color{red!60!black}$\times$ Pendiente} \\
    \midrule
    {\scriptsize 6} &
    {\tiny Análisis Socio-espacial — Diagnóstico territorial e implicaciones urbanas} &
    {\tiny\color{red!60!black}$\times$ Pendiente} \\[0.02cm]
    {\scriptsize 7} &
    {\tiny Síntesis, Lineamientos y Propuesta de Intervención} &
    {\tiny\color{red!60!black}$\times$ Pendiente} \\
    \bottomrule
  \end{tabularx}}
  \vspace{0.10cm}
  {\tiny\itshape\color{gray!55}
    {\color{green!60!black}$\surd$}\,Completo\enspace
    {\color{orange!80!black}$\triangle$}\,Parcial\enspace
    {\color{red!60!black}$\times$}\,Pendiente de cálculo o redacción\quad·\quad
    Corte: junio 2026.}
\end{frame}

% ── Frame 5.4: Cronograma ──────────────────────────────────────────────────────
%
%  TikZ horizontal con 4 bloques y flechas.
%  Colores: verde (hecho) → verde esmeralda (en curso) → acento (próximo)
%  → gris (por venir).
%
\begin{frame}{Lo que Viene: Cronograma de la Tesis}
  \vspace{0.20cm}
  \begin{center}
  \begin{tikzpicture}[
    bloque/.style = {
      rectangle, rounded corners=4pt,
      minimum width=2.55cm, minimum height=1.15cm,
      text width=2.45cm, align=center,
      font=\scriptsize, line width=0.8pt
    },
    flecha/.style = {
      ->, >=stealth, line width=1.4pt, draw=ColorPrincipal!55
    },
  ]
    \node[bloque,
          draw=green!60!black, fill=green!7, text=green!50!black] (b1) at (0,0) {
      \textbf{Caps. 1–4}\\[0.08cm]
      {\tiny Marco teórico\\código entregado}\\[0.08cm]
      {\scriptsize\bfseries Hecho}
    };
    \node[bloque,
          draw=ColorPrincipal, fill=ColorPrincipal!7, text=ColorPrincipal] (b2) at (3.2,0) {
      \textbf{Cap. 5}\\[0.08cm]
      {\tiny Simulación completa\\Fases 2 y 3}\\[0.08cm]
      {\scriptsize\bfseries Jun–Jul 2026}
    };
    \node[bloque,
          draw=ColorAcento!80!black, fill=ColorAcento!6,
          text=ColorAcento!80!black] (b3) at (6.4,0) {
      \textbf{Cap. 6}\\[0.08cm]
      {\tiny Análisis\\Socio-espacial}\\[0.08cm]
      {\scriptsize\bfseries Sept. 2026}
    };
    \node[bloque,
          draw=gray!50, fill=gray!5, text=gray!55] (b4) at (9.6,0) {
      \textbf{Cap. 7}\\[0.08cm]
      {\tiny Síntesis\\Versión Sínodo}\\[0.08cm]
      {\scriptsize\bfseries Dic. 2026}
    };
    \draw[flecha] (b1.east) -- (b2.west);
    \draw[flecha] (b2.east) -- (b3.west);
    \draw[flecha] (b3.east) -- (b4.west);
  \end{tikzpicture}
  \end{center}
  \vspace{0.20cm}
  {\scriptsize\itshape\color{ColorPrincipal}
    La reproducibilidad está garantizada: el repositorio \texttt{zmvm-transport-gis}
    contiene el código, los datos GTFS y los entornos de ejecución
    — cualquier resultado es verificable y replicable.}
\end{frame}

% ── Frame 5.5: Cierre [plain] ─────────────────────────────────────────────────
%
%  Diapositiva sin encabezado ni pie. Centrado vertical.
%  Frase de cierre entre dos líneas doradas.
%  Nombre, programa, tutor y fecha en texto pequeño al final.
%
\begin{frame}[plain]
  \vfill
  \begin{center}
    {\color{ColorAcento}\rule{0.65\textwidth}{1.2pt}}\\[0.42cm]
    {\large\itshape\color{ColorPrincipal}
      «El Periférico ya existe.\\[0.10cm]
      La pregunta no es si es posible —\\[0.06cm]
      es si la ciudad puede darse el lujo\\
      de no construirlo.»}\\[0.42cm]
    {\color{ColorAcento}\rule{0.65\textwidth}{1.2pt}}\\[0.36cm]
    {\normalsize\bfseries\color{ColorPrincipal} \Autor}\\[0.10cm]
    {\small\color{gray!65} \ProgramaCorto}\\[0.05cm]
    {\small\color{gray!65} Tutor: \ProfesorAsignado}\\[0.05cm]
    {\small\color{gray!65} \FechaPresent}
  \end{center}
  \vfill
\end{frame}
